\section{慢变量：家庭、文书与制度的尺度}
\label{sec:society-material}

\subsection{县廷把国家做成可登记的事务}

三国的国家能力并不只出现在洛阳、成都和建业的军令中。走马楼吴简使一个县廷
里的户名、田租、吏役和勾改变得可见：税役需要书手、保管、地方中介和实际
行走的人，才会从一条登记变成粮食或劳力。正因为材料出自临湘县的残存文书，
它照亮的是局部的制度实践，不能外推为全吴人口、全境亩数或一套无差别的税制。

家庭也不是财政的被动容器。一个人被写入簿籍，意味着他或她可能被纳入役使、
供给与追责；未被保存下来的名字不意味着没有劳动和关系。战争中的转输由
县廷和地方网络承受，这使北伐、江防和西晋接收旧地都具有一条很慢的社会尾声。

\subsection{名额、钱币与役使的去向}

黄初元年陈群所议的\eventanchor{TK-DENS-0220-01}{九品官人法}，让资望品评
成为魏晋选官的一种方式。它和县廷簿籍都把人转化为可处理的类别，却不可混为
同一套制度：前者的条文和纪传不能证明地方如何实际取人，后者的临湘简也不能
反推洛阳的选官网络。二三五年洛阳持续营建宫殿园囿，
\eventanchor{TK-DENS-0235-01}{役使之争把中枢工程和劳力成本显露出来}；没有
连续账簿，不能由此计算徭役总量。

孙吴对流通的压力则留下另一种痕迹。赤乌元年孙权
\eventanchor{TK-DENS-0236-01}{铸行大钱}，高额面值或许意在应对军政支出和流通，
但实施范围、兑换关系和购买力无从由官修记载确定。同年后期校事吕壹被处死，
\eventanchor{TK-DENS-0238-01}{以案件和告讦介入地方官吏的活动告终}。传记所见
是近臣、监察与地方官的张力，不能把控罪逐项当作全吴行政的统计表。

\subsection{水面上的账与损失}

黄龙二年卫温、诸葛直远航归还后受罚，
\eventanchor{TK-DENS-0230-01}{海上探求没有形成稳定据点}。记载保留了远航、
归还和处分，却不能精确对应夷洲、亶洲的地点，也不能确认死亡人数。与县廷
简牍一样，它的价值在于让国家必须负担的船、人和预期显影，而不在于替读者
画出一条已经掌握的海上疆界。

\subsection{石经、宗教与可见性的差别}

\eventanchor{TK-0240-ZHENGSHI-STONE-CLASSICS}{正始三体石经}的残石让我们看见
洛阳太学的校读与字形政治。它说明经典可以被放到公共、官学的空间中，不能
证明所有地区都以同样的文本或书体生活。宗教文本、碑刻和墓葬保存的则是另一
组被选择过的声音：成书、抄写、纪念与所述事件的日期必须区分，神异叙事也
不应被误作可直接核对的行政档案。

\subsection{律文与政权转换的不同速度}

\eventanchor{TK-0268-TAISHI-CODE}{《泰始律》的颁行}给西晋一个整理刑法的共同
语言。法典可以规定分类和程序，却不能自动缩短官吏、案件和地方之间的距离。
同样，二八〇年的统一会先改变诏令、年号和官名，随后才在州郡、港口与县廷中
遭遇不同的执行速度。制度史应当把文本当作可援引的规则，而非已经完成的社会
现实。

这种不等速也写在伐吴的物质准备里。五丈原相持时诸葛亮
\eventanchor{TK-DENS-0234-01}{分兵屯田}，只能说明军队曾试图以近处生产支撑
前线，不能倒推粮额或全部屯田范围。西陵战后，晋廷在二七三年整顿荆州军粮和
舟师，二七四年考察益、荆水陆路线，
\eventanchor{TK-DENS-0273-01}{持续的供应而非一纸诏令维持压力}；
\eventanchor{TK-DENS-0274-01}{路线考察}说明上游与荆州必须同时进入筹算。
王濬二七五年造船练水军，二七七年又积蓄舟师、军粮，朝议仍有谨慎声音：
\eventanchor{TK-DENS-0275-01}{船材、船工和训练把统一准备拖入多年}，
\eventanchor{TK-DENS-0277-01}{军粮也没有脱离争论}。史料足以显示准备的次序，
无法还原每次会议或造船规模。

\begin{turningpoint}[材料的尺度]
一枚简、一块残石和一部律文分别见证地方事务、官学文本与中央规范。把它们
并置能够抵抗只从帝王与大战写三国的习惯；把它们混同，则会制造一张并不存在
的全境总表。
\end{turningpoint}
